\documentclass[preview]{standalone}

\usepackage{amsmath}
\usepackage{amssymb}
\usepackage{stellar}
\usepackage{definitions}

\begin{document}

\id{modules-over-a-ring}
\genpage

\section{Modules over a ring}

\begin{snippetdefinition}{module-definition}{Module over a ring}
    Let \(A\) be a \ring with identity. A \emph{left \(A\)-module} is an
    \abeliangroup \((M,+)\) together with a \function
    \[
        A\times M\fromto M,
        \qquad
        (a,v)\longmapsto av,
    \]
    such that, for all \(a,b\in A\) and \(v,v_1,v_2\in M\),
    \[
        1_Av=v,
        \qquad
        (ab)v=a(bv),
    \]
    \[
        (a+b)v=av+bv,
        \qquad
        a(v_1+v_2)=av_1+av_2.
    \]
    If \(A\) is commutative, one simply says \emph{\(A\)-module}.
\end{snippetdefinition}

\begin{snippet}{modules-generalize-vector-spaces}
    A \vectorspace over a \field \(K\) is precisely a \(K\)-\module.
    Modules allow scalars in an arbitrary \ring, so nonzero scalars need
    not be invertible.
\end{snippet}

\begin{snippetproposition}{module-elementary-properties}{Elementary properties of modules}
    Let \(A\) be a \ring and let \(M\) be a left \(A\)-\module. For every
    \(a\in A\) and \(v\in M\),
    \[
        0_Av=0_M,
        \qquad
        a0_M=0_M,
    \]
    and
    \[
        a(-v)=-(av),
        \qquad
        (-a)v=-(av).
    \]
\end{snippetproposition}

\begin{snippetproof}{module-elementary-properties-proof}{module-elementary-properties}{Elementary properties of modules}
    Since
    \[
        0_Av=(0_A+0_A)v=0_Av+0_Av,
    \]
    cancellation in the additive group gives \(0_Av=0_M\). Similarly,
    \[
        a0_M=a(0_M+0_M)=a0_M+a0_M
    \]
    gives \(a0_M=0_M\).

    Finally,
    \[
        av+a(-v)=a(v-v)=a0_M=0_M
    \]
    and
    \[
        av+(-a)v=(a-a)v=0_Av=0_M.
    \]
    These identities give the two formulas for additive inverses.
\end{snippetproof}

\section{Examples}

\begin{snippetexample}{vector-spaces-as-modules-example}{Vector spaces}
    Let \(K\) be a \field. Every \vectorspace over \(K\) is a
    \(K\)-\module, with its usual scalar multiplication.
\end{snippetexample}

\begin{snippetexample}{polynomial-ring-module-example}{Polynomial rings as modules}
    Let \(A\) be a \commutativering. The polynomial \ring \(A[x]\) is an
    \(A\)-\module under
    \[
        a\cdot f=af,
        \qquad
        a\in A,
        \quad
        f\in A[x].
    \]
    Here only the additive group and scalar multiplication are retained;
    the multiplication of two polynomials is additional ring structure.
\end{snippetexample}

\begin{snippetexample}{matrix-ring-natural-module-example}{The natural module of a matrix ring}
    Let \(K\) be a \field and let
    \[
        A=\matrices_n(K).
    \]
    Matrix-vector multiplication makes \(K^n\) a left \(A\)-\module. This
    example is one reason to allow the scalar ring in the module
    definition to be noncommutative.
\end{snippetexample}

\begin{snippetexample}{ideals-and-quotients-as-modules-example}{Ideals and quotients as modules}
    Let \(A\) be a \ring. Every left \ideal \(I\subseteq A\) is a left
    \(A\)-\module under multiplication:
    \[
        a\cdot x=ax,
        \qquad
        a\in A,
        \quad
        x\in I.
    \]
    In particular, \(A\) is a left \module over itself.

    If \(I\idealin A\) is a two-sided ideal, then \(A/I\) is a left
    \(A\)-\module under
    \[
        a\cdot(b+I)=ab+I.
    \]
    This action is well-defined because \(I\) is an \ideal.
\end{snippetexample}

\end{document}
